% =============================================================
%  System Description — Adaptive and Continual Learning for HAR
%  Ready to \input{} into your main thesis document
% =============================================================

\chapter{Proposed System}

\section{Problem Formulation}

Let $\mathcal{D} = \{(\mathbf{x}_i, y_i, s_i)\}_{i=1}^{N}$ be a dataset of
inertial signals, where $\mathbf{x}_i \in \mathbb{R}^{T \times C}$ is a
time-series window of $T$ timesteps and $C$ sensor channels,
$y_i \in \mathcal{Y}$ is the activity label, and $s_i$ is the subject
identifier.

We address two complementary objectives:

\begin{enumerate}
    \item \textbf{Continual Activity Recognition.}
    Given a sequence of tasks $\{\mathcal{T}_1, \mathcal{T}_2, \ldots,
    \mathcal{T}_K\}$ arriving over time — each representing a new user,
    a new activity class, or a new context — train a model $f_\theta$ that
    performs well on \emph{all} tasks seen so far without retraining from
    scratch and without forgetting previously acquired knowledge
    (\emph{catastrophic forgetting}).

    \item \textbf{Activity Anticipation.}
    Given a partial observation $\mathbf{x}^{(p)} \in \mathbb{R}^{pT \times C}$
    ($p \in (0,1)$) of an ongoing activity, predict the label $\hat{y}^{+1}$
    of the \emph{next} activity before it begins.
\end{enumerate}


% ---------------------------------------------------------------
\section{Data Preprocessing Pipeline}
% ---------------------------------------------------------------

Raw inertial signals collected from heterogeneous datasets differ in sampling
frequency, unit of measurement, and presence of gravitational acceleration.
We apply a three-step homogenisation procedure inspired by
\cite{amrani2025leveraging}.

\subsection{Signal Homogenisation}

\paragraph{Resampling.}
All signals are resampled to a common frequency $f_s = 50$\,Hz using
Fourier-based resampling. The number of output samples is:
\begin{equation}
    N_{\text{new}} = \left\lfloor N_{\text{orig}}
        \times \frac{f_s}{f_{\text{orig}}} \right\rfloor
    \label{eq:resample}
\end{equation}

\paragraph{Unit Conversion.}
Accelerometer channels expressed in gravitational units ($g$) are converted
to SI units:
\begin{equation}
    \mathbf{a}_{[\text{m/s}^2]} = \mathbf{a}_{[g]} \times g_0,
    \qquad g_0 = 9.80665\;\text{m/s}^2
    \label{eq:unit}
\end{equation}

\paragraph{Gravity Removal.}
Gravity is estimated via a 5\textsuperscript{th}-order Butterworth low-pass
filter with cutoff $f_c = 0.5$\,Hz and removed by subtraction:
\begin{equation}
    \tilde{\mathbf{a}}(t) = \mathbf{a}(t) - \mathcal{F}_{\text{LP}}(\mathbf{a})(t)
    \label{eq:gravity}
\end{equation}
where $\mathcal{F}_{\text{LP}}$ denotes the low-pass filter operator.
This step is applied uniformly regardless of whether gravity is declared
present in the dataset metadata, following \cite{amrani2025leveraging}.

\subsection{Label Homogenisation}

Each dataset uses its own activity vocabulary. We define a unified label
space $\mathcal{Y}_u$ of 12 activity classes (Table~\ref{tab:labels}) and
map each dataset-specific label to the closest semantic equivalent using a
combination of syntactic string similarity (Levenshtein distance) and
signal-level similarity (Euclidean distance between hand-crafted feature
vectors).

\begin{table}[h]
\centering
\caption{Unified activity label space $\mathcal{Y}_u$.}
\label{tab:labels}
\begin{tabular}{clll}
\hline
\textbf{ID} & \textbf{Activity} & \textbf{Datasets} & \textbf{Windows} \\
\hline
1  & Walking              & HAPT, WISDM, PAMAP2 & 17\,349 \\
2  & Walking Upstairs     & HAPT, PAMAP2        &  6\,374 \\
3  & Walking Downstairs   & HAPT, PAMAP2        &  5\,477 \\
4  & Sitting / Standing   & HAPT, WISDM, PAMAP2 &  9\,593 \\
5  & Jogging / Running    & WISDM, PAMAP2       & 12\,048 \\
6  & Lying                & HAPT, PAMAP2        &  3\,112 \\
7  & Cycling              & PAMAP2              &  1\,095 \\
8  & Nordic Walking       & PAMAP2              &  1\,255 \\
9  & Stand $\to$ Sit      & HAPT               &    138 \\
10 & Sit $\to$ Stand      & HAPT               &    109 \\
11 & Sit $\to$ Lie        & HAPT               &    163 \\
12 & Lie $\to$ Sit        & HAPT               &    148 \\
\hline
\end{tabular}
\end{table}

\subsection{Sliding Window Segmentation}

Each preprocessed signal is divided into overlapping windows of fixed length
$T = f_s \times W_{\text{sec}} = 50 \times 3 = 150$ samples with a stride
of $S = T \times (1 - r_{\text{overlap}}) = 75$ samples
($r_{\text{overlap}} = 0.5$).
The label of each window is determined by majority vote over its sample-level
labels:
\begin{equation}
    y_w = \arg\max_{c \in \mathcal{Y}_u}
          \sum_{t=t_0}^{t_0+T-1} \mathbf{1}[y(t) = c]
    \label{eq:window_label}
\end{equation}
This yields $N = 56{,}861$ labelled windows across 44 subjects from three
publicly available datasets (HAPT, WISDM, PAMAP2).


% ---------------------------------------------------------------
\section{Model Architecture}
% ---------------------------------------------------------------

The proposed system consists of a shared \emph{backbone} encoder and two
task-specific heads, as illustrated in Figure~\ref{fig:architecture}.

\begin{figure}[h]
\centering
% Replace with actual figure when available
\fbox{\parbox{0.85\linewidth}{\centering\vspace{1.5cm}
\texttt{[Architecture diagram: IMU → Backbone → Module 1 / Module 2]}
\vspace{1.5cm}}}
\caption{Overview of the proposed two-headed continual HAR architecture.}
\label{fig:architecture}
\end{figure}

\subsection{IMU Transformer Encoder (Backbone)}

Inspired by \cite{dirgova2022transformer}, we adopt a Transformer encoder
as the shared feature extractor. The backbone maps each window
$\mathbf{x} \in \mathbb{R}^{T \times C}$ to a fixed-length embedding
vector $\mathbf{e} \in \mathbb{R}^{d}$.

\paragraph{Input Projection.}
A linear layer followed by Layer Normalisation projects the raw $C$-channel
input to the model dimension $d$:
\begin{equation}
    \mathbf{H}^{(0)} = \text{LayerNorm}(\mathbf{x}\,\mathbf{W}_{\text{in}} + \mathbf{b}_{\text{in}}),
    \quad \mathbf{W}_{\text{in}} \in \mathbb{R}^{C \times d}
\end{equation}

\paragraph{CLS Token.}
A learnable class token $\mathbf{v}_{\text{cls}} \in \mathbb{R}^d$ is
prepended to the projected sequence:
\begin{equation}
    \tilde{\mathbf{H}}^{(0)} =
    [\mathbf{v}_{\text{cls}}\;;\;\mathbf{H}^{(0)}]
    \in \mathbb{R}^{(T+1) \times d}
\end{equation}

\paragraph{Positional Encoding.}
A combination of fixed sinusoidal and learnable positional encodings is
added:
\begin{equation}
    \mathbf{H} = \tilde{\mathbf{H}}^{(0)} + \lambda \cdot \mathbf{P},
    \qquad
    P_{t,2k}   = \sin\!\left(\frac{t}{10000^{2k/d}}\right),\quad
    P_{t,2k+1} = \cos\!\left(\frac{t}{10000^{2k/d}}\right)
\end{equation}
where $\lambda \in \mathbb{R}$ is a learnable scaling parameter initialised
to 1.

\paragraph{Transformer Encoder Blocks.}
The sequence is processed by $L = 4$ identical encoder blocks.
Each block applies pre-normalisation multi-head self-attention followed
by a position-wise feed-forward network:
\begin{align}
    \mathbf{Z}^{(\ell)}   &= \mathbf{H}^{(\ell-1)}
        + \text{MHA}\!\left(\text{LN}(\mathbf{H}^{(\ell-1)})\right) \\
    \mathbf{H}^{(\ell)}   &= \mathbf{Z}^{(\ell)}
        + \text{FFN}\!\left(\text{LN}(\mathbf{Z}^{(\ell)})\right)
\end{align}
Multi-head attention with $H = 4$ heads computes:
\begin{equation}
    \text{MHA}(\mathbf{Q},\mathbf{K},\mathbf{V}) =
    \text{Concat}(\text{head}_1, \ldots, \text{head}_H)\,\mathbf{W}^O,
    \qquad
    \text{head}_h = \text{softmax}\!\left(
        \frac{\mathbf{Q}\mathbf{W}^Q_h\,(\mathbf{K}\mathbf{W}^K_h)^\top}
             {\sqrt{d/H}}
    \right)\mathbf{V}\mathbf{W}^V_h
\end{equation}
The feed-forward sub-layer consists of two linear transformations with
a GELU activation and dimension $d_{\text{ff}} = 2d = 256$:
\begin{equation}
    \text{FFN}(\mathbf{z}) =
    \text{GELU}(\mathbf{z}\,\mathbf{W}_1 + \mathbf{b}_1)\,\mathbf{W}_2
    + \mathbf{b}_2
\end{equation}

\paragraph{Output.}
The final embedding is the CLS token of the last block after Layer
Normalisation:
\begin{equation}
    \mathbf{e} = \text{LN}(\mathbf{H}^{(L)})[0] \in \mathbb{R}^d
    \label{eq:embedding}
\end{equation}
The backbone has $531{,}457$ trainable parameters.


% ---------------------------------------------------------------
\subsection{Module 1 — Continual HAR Head}
% ---------------------------------------------------------------

\subsubsection{Prototype Memory}

We adopt a \emph{prototype-based} nearest-mean classifier inspired by
LAPNet-HAR \cite{adaimi2022lifelong}. Each activity class $c \in \mathcal{Y}$
is represented by a single prototype vector $\boldsymbol{\mu}_c \in
\mathbb{R}^d$, defined as the running mean of all embeddings observed for
that class:
\begin{equation}
    \boldsymbol{\mu}_c^{(n+m)} =
    \frac{n\,\boldsymbol{\mu}_c^{(n)} + m\,\bar{\mathbf{e}}_c^{\text{new}}}
         {n + m}
    \label{eq:proto_update}
\end{equation}
where $n$ and $m$ are the previous and new sample counts for class $c$,
and $\bar{\mathbf{e}}_c^{\text{new}}$ is the mean embedding of the new batch.
This update is performed online without storing any past embeddings,
making it memory-efficient.

\paragraph{Nearest-Mean Classification.}
At inference, a window $\mathbf{x}$ is classified by finding the closest
prototype in Euclidean space:
\begin{equation}
    \hat{y} = \arg\min_{c \in \mathcal{K}}
    \left\|\mathbf{e} - \boldsymbol{\mu}_c\right\|_2^2
    \label{eq:nmc}
\end{equation}
where $\mathcal{K}$ is the set of all classes seen so far. This
\emph{fixed-capacity} head naturally accommodates new classes without
modifying the network architecture.

\subsubsection{Experience Replay Buffer}

To mitigate catastrophic forgetting, we maintain a fixed-size replay
buffer $\mathcal{M}$ of capacity $|\mathcal{M}| = 2{,}000$ windows using
class-balanced reservoir sampling. The buffer is partitioned evenly across
all known classes: each class retains at most
$\lfloor|\mathcal{M}|/|\mathcal{K}|\rfloor$ samples.

When a new sample $(\mathbf{x}, y)$ arrives for class $c$, reservoir
sampling decides whether it replaces an existing entry:
\begin{equation}
    \text{replace with probability } \frac{|\mathcal{M}|/|\mathcal{K}|}{n_c + 1}
\end{equation}
where $n_c$ is the total number of samples seen for class $c$ so far.

During continual training, a replay batch $\mathcal{B}_r \sim \mathcal{M}$
is mixed with the current task batch $\mathcal{B}_t$, exposing the model
to past distributions at every update step.

\subsubsection{Supervised Contrastive Loss}

To enforce inter-class separation in the embedding space as new classes
arrive, we employ the supervised contrastive loss of
\cite{khosla2020supervised}. Given a batch of $B$ normalised embeddings
$\hat{\mathbf{e}}_i = \mathbf{e}_i / \|\mathbf{e}_i\|_2$:
\begin{equation}
    \mathcal{L}_{\text{contrast}} =
    -\frac{1}{B} \sum_{i=1}^{B}
    \frac{1}{|\mathcal{P}(i)|}
    \sum_{j \in \mathcal{P}(i)}
    \log
    \frac{\exp\!\left(\hat{\mathbf{e}}_i^\top \hat{\mathbf{e}}_j / \tau\right)}
    {\displaystyle\sum_{k \neq i} \exp\!\left(\hat{\mathbf{e}}_i^\top \hat{\mathbf{e}}_k / \tau\right)}
    \label{eq:contrastive}
\end{equation}
where $\mathcal{P}(i) = \{j \neq i : y_j = y_i\}$ is the set of positive
pairs for anchor $i$ and $\tau = 0.07$ is the temperature.

\subsubsection{Continual Training Objective}

The total loss at each continual update step combines three terms:
\begin{equation}
    \mathcal{L}_{\text{total}} =
    \underbrace{\mathcal{L}_{\text{CE}}(\mathcal{B}_t)}_{\text{current task}}
    +\; \underbrace{\mathcal{L}_{\text{CE}}(\mathcal{B}_r)}_{\text{replay}}
    +\; \lambda_c\;\underbrace{\mathcal{L}_{\text{contrast}}(\mathcal{B}_t)}_{\text{inter-class separation}}
    \label{eq:total_loss}
\end{equation}
where $\mathcal{L}_{\text{CE}}$ is the cross-entropy loss on the linear
classification head and $\lambda_c = 0.1$ controls the contrastive
regularisation strength.

After each task, prototype vectors are \emph{re-adapted} using the replay
buffer to account for the backbone's updated embedding space:
\begin{equation}
    \boldsymbol{\mu}_c \leftarrow
    \frac{1}{|\mathcal{M}_c|}
    \sum_{\mathbf{x} \in \mathcal{M}_c} f_\theta(\mathbf{x})
    \quad \forall c \in \mathcal{K}
    \label{eq:proto_adapt}
\end{equation}
This prevents \emph{prototype obsolescence} — the misalignment between
stored prototypes and the current embedding space after weight updates
\cite{adaimi2022lifelong}.


% ---------------------------------------------------------------
\subsection{Module 2 — Activity Anticipation Head}
% ---------------------------------------------------------------

The anticipation head predicts the label of the \emph{next} activity from
a partial observation of the current one. Given a sequence of $S$ partial
windows $(\mathbf{x}^{(1)}, \ldots, \mathbf{x}^{(S)})$, each of length
$pT$ ($p < 1$), the backbone produces a sequence of embeddings:
\begin{equation}
    \mathbf{e}^{(k)} = f_\theta(\mathbf{x}^{(k)}) \in \mathbb{R}^d,
    \quad k = 1, \ldots, S
\end{equation}

These are fed into a 2-layer LSTM with hidden dimension $d_h = 128$ that
models the temporal progression of the activity stream:
\begin{equation}
    \mathbf{h}^{(k)}, \mathbf{c}^{(k)} =
    \text{LSTM}\!\left(\mathbf{e}^{(k)}, \mathbf{h}^{(k-1)},
    \mathbf{c}^{(k-1)}\right)
\end{equation}

The final hidden state is passed through a linear classifier to produce
next-activity logits:
\begin{equation}
    \hat{\mathbf{y}}^{+1} =
    \text{softmax}\!\left(
        \text{LN}(\mathbf{h}^{(S)})\,\mathbf{W}_{\text{out}}
    \right) \in \mathbb{R}^{|\mathcal{Y}|}
    \label{eq:anticipation}
\end{equation}

The model is trained to anticipate at three observation ratios
$p \in \{0.25, 0.50, 0.75\}$, encouraging robustness to partial inputs
of varying completeness. The training loss is cross-entropy between
$\hat{\mathbf{y}}^{+1}$ and the true next-window label.


% ---------------------------------------------------------------
\section{Training Procedure}
% ---------------------------------------------------------------

Training proceeds in two stages.

\subsection{Stage 1 — Supervised Pre-training}

The backbone and a linear softmax head are trained jointly on the full
merged dataset using standard cross-entropy loss, with the AdamW
optimiser \cite{loshchilov2019adamw} ($\beta_1{=}0.9$, $\beta_2{=}0.999$,
weight decay $= 10^{-4}$) and a cosine annealing learning rate schedule:
\begin{equation}
    \eta_t = \eta_{\min} + \frac{1}{2}(\eta_0 - \eta_{\min})
    \left(1 + \cos\frac{\pi\,t}{T_{\max}}\right)
    \label{eq:cosine_lr}
\end{equation}
with $\eta_0 = 10^{-3}$, $\eta_{\min} = 10^{-5}$, $T_{\max} = 30$ epochs,
batch size 128. Gradient norms are clipped at 1.0.

At the end of pre-training, prototype vectors are initialised from the
training set embeddings using Eq.~\eqref{eq:proto_update}.

\subsection{Stage 2 — Continual Learning}

Tasks arrive sequentially. For each task $\mathcal{T}_i$, the model is
updated for $E = 10$ epochs using the composite loss in
Eq.~\eqref{eq:total_loss} with a reduced learning rate $\eta = 5 \times
10^{-5}$ to prevent catastrophic overwriting of the shared backbone.
After training on task $\mathcal{T}_i$, prototypes are re-adapted via
Eq.~\eqref{eq:proto_adapt} using the current replay buffer.

Two continual learning scenarios are evaluated:

\begin{itemize}
    \item \textbf{User-incremental:} each task $\mathcal{T}_i$ introduces
    data from a new subject $s_i$. The model must adapt to the new
    individual's movement style while retaining performance on all previous
    subjects.

    \item \textbf{Class-incremental:} each task $\mathcal{T}_i$ introduces
    two new activity classes. The model must extend its classification
    vocabulary without retraining from scratch.
\end{itemize}


% ---------------------------------------------------------------
\section{Evaluation Metrics}
% ---------------------------------------------------------------

Let $R_{i,j}$ denote the macro-F1 score on the test set of task
$\mathcal{T}_j$ evaluated \emph{after} training on task $\mathcal{T}_i$
(with $i \geq j$). The results matrix $\mathbf{R} \in \mathbb{R}^{K \times K}$
captures the full performance profile across all tasks and training stages.

\paragraph{Macro F1 Score.}
The primary per-task metric, which accounts for class imbalance:
\begin{equation}
    F1_{\text{macro}} = \frac{1}{|\mathcal{Y}|}\sum_{c \in \mathcal{Y}}
    \frac{2\,\text{TP}_c}{2\,\text{TP}_c + \text{FP}_c + \text{FN}_c}
\end{equation}

\paragraph{Backward Transfer (BWT).}
Measures whether learning new tasks degrades performance on old ones
\cite{lopez2017gradient}. A negative value indicates forgetting:
\begin{equation}
    \text{BWT} = \frac{1}{K-1}\sum_{j=1}^{K-1}
    \left(R_{K,j} - R_{j,j}\right)
    \label{eq:bwt}
\end{equation}

\paragraph{Forgetting.}
The average maximum performance drop per old task:
\begin{equation}
    \mathcal{F} = \frac{1}{K-1}\sum_{j=1}^{K-1}
    \left(\max_{i \leq K} R_{i,j} - R_{K,j}\right)
    \label{eq:forgetting}
\end{equation}

\paragraph{Forward Transfer (FWT).}
Measures whether knowledge from past tasks accelerates learning of new ones:
\begin{equation}
    \text{FWT} = \frac{1}{K-1}\sum_{i=2}^{K}
    \left(R_{i-1,i} - b_i\right)
    \label{eq:fwt}
\end{equation}
where $b_i$ is the random-chance baseline accuracy for task $\mathcal{T}_i$.

\paragraph{Anticipation Accuracy.}
For the anticipation head, accuracy is reported at each observation ratio
$p \in \{0.25, 0.50, 0.75\}$ to characterise the trade-off between
prediction lead time and accuracy.


% ---------------------------------------------------------------
%  Bibliography entries (add to your main .bib file)
% ---------------------------------------------------------------
% @article{amrani2025leveraging,
%   title={Leveraging dataset integration and continual learning for HAR},
%   author={Amrani, Hamza and others},
%   journal={Int. J. Mach. Learn. Cyber.},
%   volume={16}, pages={5213--5234}, year={2025}
% }
% @article{adaimi2022lifelong,
%   title={Lifelong Adaptive Machine Learning for Sensor-Based HAR
%          Using Prototypical Networks},
%   author={Adaimi, Rebecca and Thomaz, Edison},
%   journal={Sensors}, volume={22}, number={18}, pages={6881}, year={2022}
% }
% @inproceedings{khosla2020supervised,
%   title={Supervised Contrastive Learning},
%   author={Khosla, Prannay and others},
%   booktitle={NeurIPS}, year={2020}
% }
% @article{dirgova2022transformer,
%   title={Wearable Sensor-Based HAR with Transformer Model},
%   author={Dirgov{\'a} Lupt{\'a}kov{\'a}, Ivana and others},
%   journal={Sensors}, volume={22}, number={5}, pages={1911}, year={2022}
% }
% @inproceedings{lopez2017gradient,
%   title={Gradient Episodic Memory for Continual Learning},
%   author={Lopez-Paz, David and Ranzato, Marc'Aurelio},
%   booktitle={NeurIPS}, year={2017}
% }
% @inproceedings{loshchilov2019adamw,
%   title={Decoupled Weight Decay Regularization},
%   author={Loshchilov, Ilya and Hutter, Frank},
%   booktitle={ICLR}, year={2019}
% }
